\documentclass[11pt]{article}

% ============================================================
% Physics Exercise Template
% ============================================================

% Mathematics
\usepackage{amsmath}
\usepackage{amssymb}
\usepackage{mathtools}
\usepackage{bm}

% Physics notation
\usepackage{physics}
\usepackage{siunitx}
\usepackage{braket}

% Typography and layout
\usepackage{microtype}
\usepackage[margin=1in]{geometry}
\usepackage{enumitem}

% Figures and references
\usepackage{graphicx}
\usepackage{float}
\usepackage[hidelinks]{hyperref}
\usepackage{cleveref}

% ============================================================
% Basic notation
% ============================================================

\newcommand{\R}{\mathbb{R}}
\newcommand{\C}{\mathbb{C}}
\newcommand{\N}{\mathbb{N}}
\newcommand{\Z}{\mathbb{Z}}

\newcommand{\id}{\mathbb{1}}
\newcommand{\diag}{\operatorname{diag}}

% ============================================================
% Document information
% ============================================================

\newcommand{\studentname}{Heewon Lee}
\newcommand{\coursecode}{PH.60053}
\newcommand{\coursename}{Relativistic Quantum Field Theory I}
\newcommand{\semester}{Fall 2026}


\title{\coursename (\coursecode)\\Problem Set 1}
\author{\studentname}
\date{\semester}

\begin{document}

\maketitle

\paragraph{Note on units:}
I generally disdain the use of natural units in undergraduate studies.
For students having gone through a symbolic mess of a calculation,
something as simple as dimensional analysis gives a good sanity check.
And while this course is explicitly designed for graduate students,
I would like to stand by this conviction and perform all calculations
with explicit engineering dimensions and constants.
I believe it is reasonably simple for the grader to convert it into natural units.
I appreciate your understanding.

During class, the professor has instructed us that the engineering dimensions of $\phi$
must be determined such that $\lagrdens$ has dimensions $[E L^{-3}]$.
This means $[\phi] = [E^{1/2} L^{-1/2}]$, and the Klein-Gordon equation is given by
\[
    \left(\square + \left(\frac{mc}{\hbar}\right)^2\right)\phi
    = \left(
        -\frac{1}{c^2}\partial_t^2
        + \nablavec^2
        + \left(\frac{mc}{\hbar}\right)^2
    \right) \phi
    = 0.
\]
Suppose further that the mode expansion is given as:
\[
    \phi(t,\vec{x}) = \int\frac{d^3\vec{p}}{(2\pi\hbar)^{3/2}} \mathcal{N}(E_{\vec{p}})
    \left(
        a_{\vec{p}} e^{\frac{ip\cdot x}{\hbar}}
        + a_{\vec{p}}^\dagger e^{-\frac{ip\cdot x}{\hbar}},
    \right)
\]
where $x^\mu = (ct, \vec{x})$, $p^\mu = (E_{\vec{p}}/c, \vec{p})$, and
$E_{\vec{p}} = \sqrt{(mc^2)^2 + (\vec{p}c)^2}$.
Insisting the simple commutation relation:
\[
    \left[a_{\vec{p}}, a_{\vec{p}'}\right] = \delta^{(3)}(\vec{p}-\vec{p}')
\]
implies that the normalization factor $\mathcal{N}(E_{\vec{p}})$ has dimensions $[E^{1/2} L]$, or
\[
    d\Pi_{\vec{p}}
    := \frac{d^3\vec{p}}{(2\pi\hbar)^{3/2}} \mathcal{N}(E_{\vec{p}})
    = \frac{d^3\vec{p}\; \hbar c}{(2\pi\hbar)^{3/2}\sqrt{2E_{\vec{p}}}}.
\]

\section*{Problem 3.1}
Let us substitute the mode expansion into the definition of the Gaussian-smeared operator:
\[
    A(a)
    = \frac{1}{(a\sqrt{\pi})^3} \int d\Pi_{\vec{p}} \int d^3\vec{x} \left(
        a_{\vec{p}} e^{\frac{ip\cdot x}{\hbar}} + a_{\vec{p}}^\dagger e^{-\frac{ip\cdot x}{\hbar}}
    \right) e^{-\vec{x}^2/a^2}.
\]
Here, because the time coordinate is zero, $p\cdot x = \vec{p}\cdot\vec{x}$.
Thus, we can integrate the exponentials over the spatial coordinates first.
\begin{align*}
    \int d^3\vec{x} \exp\left(-\frac{\vec{x}^2}{a^2} + \frac{i\vec{p}\cdot\vec{x}}{\hbar}\right)
    &= \int_0^\infty dr r^2 e^{-r^2/a^2} \int_{-1}^1 d(\cos\theta) \int_0^{2\pi} d\phi
        e^{\frac{i|\vec{p}|r}{\hbar}\cos\theta} \\
    &= \int_0^\infty dr r^2 e^{-r^2/a^2} \cdot 4\pi
        \frac{\sin\frac{|\vec{p}|r}{\hbar}}{|\vec{p}|r/\hbar} \\
    &= \frac{4\pi\hbar}{|\vec{p}|} \int_0^\infty dr r\sin\frac{|\vec{p}|r}{\hbar} e^{-r^2/a^2} \\
    &= \frac{4\pi\hbar}{|\vec{p}|}
        \left(-\hbar\frac{\partial}{\partial(|\vec{p}|)}\right)
        \int_0^\infty dr \cos\frac{|\vec{p}|r}{\hbar} e^{-r^2/a^2} \\
    &= \frac{4\pi\hbar^2}{|\vec{p}|}
        \left(-\frac{\partial}{\partial(|\vec{p}|)}\right)
        \frac{1}{2} \Re \int_{-\infty}^\infty dr \exp\left(
            -\frac{r^2}{a^2} + \frac{i|\vec{p}|r}{\hbar}
        \right) \\
    &= \frac{2\pi\hbar^2}{|\vec{p}|}
        \left(-\frac{\partial}{\partial(|\vec{p}|)}\right)
        \Re \int_{-\infty}^\infty dr \exp\left(
            -\frac{1}{a^2}\left(r - \frac{i|\vec{p}|a^2}{\hbar}\right)^2
            -\frac{|\vec{p}|^2a^2}{4\hbar^2}
        \right) \\
    &= \frac{2\pi\hbar^2}{|\vec{p}|}
        \left(-\frac{\partial}{\partial(|\vec{p}|)}\right)
        \sqrt{\pi}a \cdot \exp\left(-\frac{|\vec{p}|^2 a^2}{4\hbar^2}\right) \\
    &= (a\sqrt{\pi})^3 \exp\left(-\frac{|\vec{p}|^2 a^2}{4\hbar^2}\right)
\end{align*}
\[
    \Rightarrow A(a)
    = \int d\Pi_{\vec{p}} \left( a_{\vec{p}} + a_{\vec{p}}^\dagger \right)
        \exp\left(-\frac{|\vec{p}|^2 a^2}{4\hbar^2}\right)
\]

Since $a_{\vec{p}}\ket{0} = \bra{0} a_{\vec{p}}^\dagger = 0$,
this shows that $\braket{0 | A(a) | 0} = 0$.
\begin{align*}
    \Rightarrow \var A(a)
    &= \braket{0 | A(a)^2 | 0} \\
    &= \int d\Pi_{\vec{p}} \int d\Pi_{\vec{p}'}
        \braket{0 | \left(
            a_{\vec{p}}a_{\vec{p}'}
            + a_{\vec{p}}a_{\vec{p}'}^\dagger
            + a_{\vec{p}}^\dagger a_{\vec{p}'}
            + a_{\vec{p}}^\dagger a_{\vec{p}'}^\dagger
        \right) | 0} \exp \left(
            -\frac{(|\vec{p}|^2 + |\vec{p}'|^2)a^2}{4\hbar^2}
        \right) \\
    &= \int d\Pi_{\vec{p}} \int d\Pi_{\vec{p}'}
        \braket{0 | a_{\vec{p}}a_{\vec{p}'}^\dagger | 0} \exp \left(
            -\frac{(|\vec{p}|^2 + |\vec{p}'|^2)a^2}{4\hbar^2}
        \right) \\
    &= \int d\Pi_{\vec{p}} \int d\Pi_{\vec{p}'}
        \braket{0 | \left(
            a_{\vec{p}'}^\dagger a_{\vec{p}}
            + \delta^{(3)}(\vec{p} - \vec{p}')
        \right) | 0} \exp \left(
            -\frac{(|\vec{p}|^2 + |\vec{p}'|^2)a^2}{4\hbar^2}
        \right) \\
    &= \int\frac{d^3\vec{p}\; \hbar^2 c^2}{(2\pi\hbar)^3 \cdot 2E_{\vec{p}}}
        \exp\left(-\frac{|\vec{p}|^2 a^2}{2\hbar^2}\right) \\
    &= \int\frac{
        d^3\vec{\kappa}\; \hbar^2 c^2 (\hbar/a)^3
    }{
        (2\pi\hbar)^3 \cdot 2\hbar c / a \; \sqrt{\vec{\kappa}^2 + (mca/\hbar)^2}
    } e^{-\vec{\kappa}^2/2}
            \qquad (\vec{\kappa} := \vec{p}a / \hbar) \\
    &= \frac{\hbar c}{16\pi^3 a^2} \int d^3\vec{\kappa}
        \frac{e^{-\vec{\kappa}^2/2}}{\sqrt{\vec{\kappa}^2 + (mca/\hbar)^2}} \\
    &= \frac{\hbar c}{4\pi^2 a^2} \int_0^\infty d\kappa
        \frac{\kappa^2 e^{-\kappa^2/2}}{\sqrt{\kappa^2 + (mca/\hbar)^2}}
\end{align*}

Let us define the dimensionless parameter $\mu := mca/\hbar$ and find the asymptotic behavior:
\begin{itemize}
    \item $\mu \gg 1$:
        \begin{align*}
            \var A(a)
            &= \frac{\hbar c}{4\pi^2 a^2} \int d\kappa
                \frac{\kappa^2}{\mu} e^{-\kappa^2/2} \left(
                    1 + \left(\frac{\kappa}{\mu}\right)^2
                \right)^{-1/2} \\
            &= \frac{\hbar c}{4\pi^2 a^2 \mu} \int d\kappa
                e^{-\kappa^2/2}\left(\kappa^2 - \frac{\kappa^4}{2\mu^2} + \cdots\right) \\
            &\sim \frac{\hbar c}{\sqrt{32\pi^3}a^2 \mu} \left(1 + O(\mu^{-2})\right) \\
            &= \left(\frac{\hbar^2}{\sqrt{32\pi^3} m}\right) a^{-3} \cdot \left(
                1 + O\left(\left(\frac{mca}{\hbar}\right)^{-2}\right)
            \right)
        \end{align*}
        Smearing over regions much larger than the correlation length of the field,
        the average fluctuation inside the region scales like the volume of the region.
        This can in understood in terms of field correlations decaying to zero at long distances,
        rendering the fluctuations essentially independent between separated regions.

    \item $\mu \ll 1$:
        \begin{align*}
            \var A(a)
            &= \frac{\hbar c}{4\pi^2 a^2} \int d\kappa
                \kappa e^{-\kappa^2/2} \left(
                    1 + \left(\frac{\mu}{\kappa}\right)^2
                \right)^{-1/2} \\
            &= \frac{\hbar c}{4\pi^2 a^2} \int d\kappa
                e^{-\kappa^2/2}\left(\kappa - \frac{\mu^2}{2\kappa} + \cdots\right) \\
            &\sim \frac{\hbar c}{4\pi^2 a^2} \left(1 + O(\mu^2)\right) \\
            &= \left(\frac{\hbar c}{4\pi^2}\right) a^{-2} \cdot \left(
                1 + O\left(\left(\frac{mca}{\hbar}\right)^2\right)
            \right)
        \end{align*}
        In regions much smaller than the correlation length,
        any fluctuation of the field inside the region must have a short wavelength,
        i.e., a large momentum of order $\hbar / a \gg mc$.
        Thus, the particles are approximately massless,
        and the field fluctuation follows a inverse-area law $a^{-2}$.
\end{itemize}

\section*{Problem 3.2}
\noindent
(a) For a complex scalar field $\phi : \mathbb{R}^{1,3} \rightarrow \mathbb{C}$,
we assert that the Lagrangian density $\lagrdens[\phi]$ must satisfy the following conditions:
\begin{itemize}
    \item \textbf{Lorentz invariance}:
        Upon any Poincar\'e transform $(a,\Lambda)$, the Lagrangian density must satisfy
        $\lagrdens[\phi]\overset{(a,\Lambda)}{\mapsto} \lagrdens'[\phi'] = \lagrdens[\phi]$.

    \item \textbf{Reality}: $\lagrdens^\dagger = \lagrdens$

    \item \textbf{Global $U(1)$ symmetry}:
        The Lagrangian density must be a function only of the ray of $\phi$,
        as the vectors in the ray correspond to the same physical state.
        That is, it must be invariant under $\phi \mapsto e^{i\alpha} \phi$.
\end{itemize}


\end{document}
